\begin{textsample}{POS Dim 1 – 2017\_c6 – Score -3.00 – 2017\_c6/enade\_2017\_letras\_por\_esp\_f454913e6e.txt}  \label{ex:f1_pos_299}
CAPÍTULO 6 ANÁLISE TÉCNICA DA PROVA Este capítulo tem por objetivo apresentar o desempenho dos estudantes concluintes de Letras-Português e Espanhol ( Licenciatura ) no Enade/2017.
Para isso, foram calculadas as estatísticas básicas da prova em seu todo ( seção 6.1.1 ), bem como as estatísticas relacionadas aos Componentes de Formação Geral ( seção 6.1.2 ) e de Conhecimento Específico da Área ( seção 6.1.3 ).
Dadas as suas características, foram analisadas, em separado, as questões objetivas ( seção 6.2 ) e as questões discursivas ( seção 6.3 ).
Para as questões objetivas, foram disponibilizados os índices de facilidade e de discriminação ponto-bisserial, também em separado, para os Componentes de Formação Geral ( seção 6.2.1 ) e de Conhecimento Específico ( seção 6.2.2 ).
De cada componente, uma das questões foi escolhida para exemplificar a análise gráfica, relacionando as alternativas escolhidas pelos estudantes ( inclusive o gabarito ) com o número de acertos no componente.
No Anexo I, é apresentada a íntegra da análise gráfica para todas as questões objetivas.
Para cada uma das questões discursivas, os conteúdos dos tipos mais comuns de respostas dos estudantes são apresentados e comparados com o padrão de respostas esperado ( ver Anexo VIII com o padrão de respostas ).
Tomando -se como base as duas questões discursivas do Componente de Formação Geral, na seção 6.3.1.6 são apresentados comentários sobre a correção das respostas em relação à Língua Portuguesa.
Das tabelas constam as seguintes estatísticas das notas19 : média do desempenho na prova, erro padrão da média, desvio padrão, nota mínima, mediana e nota máxima para cada um de seus componentes.
Tais estatísticas contemplam o total de estudantes concluintes da Área de Letras-Português e Espanhol ( Licenciatura ) inscritos e presentes à prova do Enade/2017, tendo -se em vista agregações, ou por Grandes Regiões e o país como um todo, ou por Categoria Administrativa, Organização Acadêmica da IES e Modalidade de Ensino.
Em relação aos gráficos de distribuição de notas, o intervalo considerado foi de 10 unidades, aberto à esquerda e fechado à direita, com exceção do primeiro intervalo, [ 0 ; 10 ], fechado em ambos os extremos.
Para os gráficos de distribuição das notas das questões discursivas, foram consideradas mais duas categorias : questão em branco20 e nota zero. 19 Uma definição dessas estatísticas pode ser encontrada no Glossário. 20 Nesse grupo estão incluídas também as respostas classificadas como nulas ou desconsideradas. 6.1 ESTATÍSTICAS BÁSICAS DA PROVA Nesta seção, apresentam -se estatísticas selecionadas e histogramas da nota geral ( 6.1.1 ) e de cada componente : Formação Geral ( 6.1.2 ) e Conhecimento Específico ( 6.1.3 ).
São também apresentadas estatísticas selecionadas de subpopulações, caracterizadas por Grande Região, Categoria Administrativa, Organização Acadêmica e Modalidade de Ensino. 6.1.1 Estatísticas Básicas Gerais Na Tabela 6.1, são apresentadas as Estatísticas Básicas da prova, por Grande Região, dos estudantes concluintes de Letras-Português e Espanhol ( Licenciatura ).
A Média das notas da prova, como um todo ( nas seções seguintes serão analisados os Componentes de Formação Geral e de Conhecimento Específico ), foi 44,2.
Os alunos das regiões Norte e Nordeste obtiveram a Média mais baixa ( 40,1 ), e os da região Sul obtiveram a Média mais alta ( 49,2 ).
As demais médias foram : 48,3 na região Sudeste e 40,3 na região Centro-Oeste.
O Desvio padrão para o Brasil, como um todo, foi 15,4, sendo o menor Desvio padrão encontrado na região Centro-Oeste ( 13,6 ) e o maior, na região Nordeste ( 15,4 ), indicando uma dispersão um pouco maior das notas da última região.
A região que obteve a maior nota Máxima foi a Sudeste ( 91,0 ), ao passo que a região que atingiu a menor nota Máxima foi a Norte ( 75,5 ).
A Mediana do Brasil, como um todo, foi 43,7, sendo a maior Mediana obtida na região Sul ( 49,7 ), e a menor obtida na região Nordeste ( 38,9 ).
A nota Mínima foi zero apenas na região Centro-Oeste.
Nas demais regiões as notas mínimas foram : 2,2 na região Norte, 4,3 na região Nordeste, 11,1 na região Sudeste e 14,2 na Sul.
Considerando -se as notas segundo Grande Região, observa -se que existe diferença estatisticamente significativa ao nível de 95 \% entre as três menores médias, obtidas nas regiões Norte, Nordeste ( 39,8 em ambas ) e Centro-Oeste ( 40,3 ), e as médias mais altas, obtidas nas regiões : Sudeste ( 48,3 ) e Sul (49,2)21.
O comportamento das notas dos estudantes de todo o Brasil pode ser observado no Gráfico 6.1, que apresenta um histograma com a sua distribuição.
Essa é uma distribuição unimodal, com a moda no intervalo ( 40 ; 50 ], seguido pelo intervalo ( 30 ; 40 ].
Na Tabela 6.2, são apresentadas informações referentes às Médias da nota final, desagregadas por Categoria Administrativa, por Organização Acadêmica e por Modalidade de Ensino.
A Média das notas da prova, como um todo, foi 44,2.
Em relação à Categoria Administrativa, os estudantes das IES Públicas obtiveram Média um pouco mais alta ( 44,9 ), e os das IES Privadas obtiveram Média mais baixa ( 42,8 ) que a Média nacional.
Observa -se que existe diferença estatisticamente significativa entre as médias das notas das IES Públicas e Privadas.
A diferença entre as médias das regiões Nordeste e Sul ( 9,1 ), a maior e a menor Média, é superior à diferença entre as médias das IES Públicas e as das Privadas ( 2,1 ), caracterizando maior diversidade regional do que administrativa.
O Desvio padrão para as IES Públicas ( 15,5 ) foi muito próximo ao do Brasil como um todo ( 15,4 ), indicando uma dispersão parecida das notas nesta Categoria Administrativa.
No tocante à Organização Acadêmica, os CEFET/IFET obtiveram a Média mais alta ( 47,1 ), e as Faculdades, a mais baixa ( 39,8 ).
A Média das Universidades foi 44,9 e a dos Centros universitários foi 44,7, mais altas do que a nacional ( 44,2 ).
Constata -se que existe diferença estatisticamente significativa ao nível de 95 \% entre a Média mais baixa, 39,8 das Faculdades, e a dos demais tipos de Organização Acadêmica.
A Média da Modalidade Educação a Distância ( 41,2 ) foi inferior à Média da Educação Presencial ( 45,2 ), e a diferença entre elas é estatisticamente significativa ao nível de 95 \%. 6.1.2 Estatísticas Básicas no Componente de Formação Geral Na Tabela 6.3, são apresentadas as Estatísticas Básicas em relação ao componente da prova que avalia a Formação Geral dos estudantes concluintes.
Os alunos de todo o Brasil obtiveram desempenho médio de 52,1.
Quanto à variabilidade, o Desvio padrão das notas dos estudantes do Brasil, como um todo, foi 17,4.
A maior Média foi obtida na região Sul ( 56,4 ), e a menor, na região Norte ( 46,5 ).
As demais médias foram : 49,2 na região Nordeste, 55,2 na região Sudeste e 49,2 na região Centro-Oeste.
Já o maior Desvio padrão foi obtido na região Nordeste ( 17,7 ), e o menor, na região Sul ( 15,6 ).
Os demais desvios padrões foram : 17,1 na região Norte, 17,3 na região Sudeste e 16,8 na região Centro-Oeste.
A maior nota no Componente de Formação Geral da prova do Enade/2017 foi 98,6, obtida por, pelo menos, um aluno na região Nordeste.
A menor nota Máxima foi obtida na região Norte ( 89,4 ).
A Mediana do Brasil, como um todo, foi 53,1, sendo a menor Mediana encontrada na região Norte ( 48,3 ), e a maior, encontrada na região Sul ( 57,2 ).
A nota Mínima nesta parte foi zero em quatro das cinco as regiões, a exceção ocorreu na região Sul, onde a nota Mínima foi 8,4.
Considerando -se as notas, segundo Grande Região, observa -se que existe diferença estatisticamente significativa das duas maiores médias das notas do Componente de Formação Geral, obtidas nas regiões Sul ( 56,4 ) e Sudeste ( 55,2 ), em relação às três menores médias, obtidas nas regiões Norte, Nordeste e Centro-Oeste.
Já a diferença entre as duas médias maiores não é estatisticamente significativa, como também entre as três menores.
No Gráfico 6.2, é possível observar a avaliação do desempenho dos estudantes do Componente de Formação Geral, a partir do histograma da distribuição das notas correspondentes.
A distribuição é unimodal, com moda em ( 50 ; 60 ], intervalo modal superior ao da distribuição de notas da prova, como um todo ( Gráfico 6.1 ).
Nota -se, ainda que, no Gráfico 6.2, as notas apresentam uma dispersão maior do que as do Gráfico 6.1 ( distribuição das notas da prova ), o que é confirmado pela comparação dos desvios padrões : 15,4 para a nota da prova, como um todo, e 17,4 para o Componente de Formação Geral.
Na Tabela 6.4, são apresentadas as informações referentes ao desempenho dos concluintes do Componente de Formação Geral, em diferentes agregações : Categoria Administrativa, Organização Acadêmica e Modalidade de Ensino.
Observa -se que existe diferença estatisticamente significativa das médias dos tipos de Categoria Administrativa : a maior Média obtida por estudantes de IES Públicas ( 52,8 ), e a menor, pelos de IES Privadas ( 50,6 ).
No caso das Organizações Acadêmicas, há diferença estatisticamente significativa apenas dentre as médias das Universidades e das Faculdades.
A maior Média foi obtida pelos estudantes dos CEFET/IFET ( 55,7 ), e, em ordem decrescente, as demais médias foram : 52,7 nas Universidades ; 49,0, nos Centros universitários ; e 48,8, nas Faculdades.
No que se refere à Modalidade de Ensino, a Educação Presencial apresentou maior Média do que a Educação a Distância ( médias 53,1 e 49,3, respectivamente ).
Tais médias apresentam diferença estatisticamente significativa. 6.1.3 Estatísticas Básicas do Componente de Conhecimento Específico Na Tabela 6.5, são apresentadas as Estatísticas Básicas referentes ao Componente de Conhecimento Específico da área de Letras-Português e Espanhol ( Licenciatura ).
A Média do desempenho dos alunos do Brasil, como um todo, foi 41,5.
A maior Média foi obtida na região Sul ( 46,7 ), e a menor, na região Nordeste ( 37,0 ).
As demais médias foram : 37,9 na região Norte, 46,0 na região Sudeste e 37,4 na região Centro-Oeste.
Quanto à variabilidade das notas, o Desvio padrão do Brasil, como um todo, foi 17,2, sendo o maior Desvio padrão observado na região Sudeste ( 17,5 ), e o menor, na região Centro-Oeste ( 15,0 ).
Os demais desvios foram : 15,3 na região Norte, 16,9 na região Nordeste e 15,7 na região Sul.
A Mediana das notas dos estudantes de todo o Brasil foi 40,7.
A maior Mediana ocorreu na região Sul ( 47,5 ), e a menor, na região Centro-Oeste ( 34,9 ).
As demais medianas foram : 37,3 na região Norte, 36,4 na região Nordeste e 45,8 na região Sudeste.
A nota Máxima do Brasil, como um todo, foi 94,7, sendo obtida por, pelo menos, um aluno na região Sudeste.
As demais notas máximas foram : 81,6 na região Norte, 93,5 na região Nordeste, 89,4 na região Sul e 82,6 na região Centro-Oeste.
A nota Mínima foi diferente de zero apenas na região Sul ( 4,1 ).
Observa -se que só existe diferença estatisticamente significativa das menores médias das notas do Componente de Conhecimento Específico obtidas nas regiões Norte ( 37,9 ), Nordeste ( 37,0 ) e Centro-Oeste ( 37,4 ) em relação às duas regiões com as médias mais altas : Sudeste ( 46,0 ) e Sul ( 46,7 ).
Assim como nos Gráficos 6.1 e 6.2, no Gráfico 6.3, apresentado a seguir, pode -se fazer uma avaliação do desempenho de concluintes de Letras-Português e Espanhol ( Licenciatura ) em relação ao Componente de Conhecimento Específico, com um histograma da distribuição das notas correspondentes.
Essa também é uma distribuição unimodal, e o grupo modal é o ( 30 ; 40 ], inferior tanto em relação ao grupo modal da prova, como um todo, quanto ao de Formação Geral.
Na Tabela 6.6, é apresentada uma comparação dos resultados em relação à Categoria Administrativa, à Organização Acadêmica e à Modalidade de Ensino, agora levando -se em conta o desempenho de alunos do Componente de Conhecimento Específico da prova.
No que se refere à Organização Acadêmica, como em Formação Geral e na prova como um todo, a maior Média em Conhecimento Específico foi obtida pelos CEFET/IFET ( 44,3 ), e a menor foi a das Faculdades ( 36,8 ).
As médias dos Centros universitários e das Universidades foram, respectivamente, 43,2 e 42,2.
O maior Desvio padrão foi o das Universidades ( 17,6 ), um pouco superior ao valor para o Brasil ( 17,2 ), como um todo.
As Universidades também obtiveram a maior nota Máxima ( 94,7 ).
As demais notas máximas foram : 81,4 nos Centros universitários, 78,3 nas Faculdades e 63,7 nos CEFET/IFET, a menor delas.
As medianas foram : 41,7 nas Universidades, 40,5 nos Centros universitários, 36,5 nas Faculdades, a menor, e 44,8 nos CEFET/IFET, a maior.
A nota Mínima foi diferente de zero nos CEFET/IFET ( 19,2 ) e nos Centros universitários ( 8,1 ).
Observa -se que existem diferenças estatisticamente significativas ao nível de 95 \% no Componente de Conhecimento Específico entre a menor Média, das Faculdades ( 36,8 ), e as médias dos demais tipos de Organização Acadêmica.
Quanto à Categoria Administrativa, o mesmo comportamento observado para a parte de Formação Geral e para a prova como um todo – existe diferença estatisticamente significativa entre as médias das IES Públicas ( 42,2 ) e as das IES Privadas ( 40,1 ).
Também, nesse caso, a maior Média foi obtida por alunos de IES Públicas de ensino.
Quanto ao comportamento das médias no que toca à Modalidade de Ensino, há diferença estatisticamente significativa entre a Média da Educação Presencial ( 42,6 ) e a da Educação a Distância ( 38,5 ). 6.2 ANÁLISE DAS QUESTÕES OBJETIVAS Nesta seção, apresentam -se estatísticas selecionadas e histogramas das Questões Objetivas de Formação Geral ( 6.2.1 ) e de Conhecimento Específico ( 6.2.2 ).
São também apresentadas e comparadas as médias das subpopulações, caracterizadas por Grande Região. 6.2.1 Componente de Formação Geral Na Tabela 6.7, apresentam -se as Estatísticas Básicas relativas às oito questões objetivas do componente da prova que abrange a Formação Geral dos estudantes.
A Média do Brasil foi 49,1.
A menor Média foi encontrada na região Norte ( 41,2 ), e a maior, na região Sul ( 53,0 ).
As demais médias foram : 46,7 na região Nordeste, 52,3 na região Sudeste e 46,4 na região Centro-Oeste.
O Desvio padrão do Brasil foi 21,0, sendo o maior Desvio padrão encontrado na região Sudeste ( 21,2 ), e o menor, na região Norte ( 19,3 ).
Os demais desvios foram : 21,0 na região Nordeste, 19,7 na região Sul e 21,1 na região Centro-Oeste.
Na Tabela 6.8, são apresentados o Índice de Facilidade e o Índice de Discriminação ( ponto-bisserial ) para cada uma das questões objetivas do Componente de Formação Geral.
Quanto ao Índice de Facilidade, foram usadas as seguintes cores para diferenciar o nível de dificuldade da questão : • Azul para as questões classificadas com índice Muito fácil ( \&gt ; =0,86 ), verde para as questões classificadas com índice Fácil ( 0,61 a 0,85 ), amarelo para as questões classificadas com Médio ( 0,41 a 0,60 ), vermelho para as questões classificadas com Difícil ( 0,16 a 0,40 ) e roxo para as questões classificadas com Muito difícil ( \&lt ; =0,15 ).
Já quanto ao Índice de Discriminação, foram usadas as seguintes cores para qualificar a questão : • As questões classificadas com índice Fraco receberam a cor vermelha ( \&lt ; =0,19 ), as classificadas com Médio receberam a cor amarela ( 0,20 a 0,29 ), as classificadas com Bom receberam a cor verde ( 0,30 a 0,39 ) e as classificadas com Muito bom ( \&gt ; =0,40 ) receberam a cor azul.
As questões objetivas do Componente de Formação Geral, segundo o Índice de Facilidade, foram assim avaliadas : das oito questões, nenhuma teve o Índice de Facilidade classificado como Muito fácil, e duas questão foram classificadas como Fácil, por terem índice de acertos entre 0,61 e 0,85.
Três questões foram consideradas com índice de dificuldade Médio, situando -se no intervalo entre 0,41 e 0,60 do Índice de Facilidade, ou seja, houve entre 41,0 \% e 60,0 \% de acertos, enquanto duas questões foram classificadas na categoria Difícil, As medianas do Brasil, como um todo e de três das cinco regiões foi 42,9.
Nas regiões Sudeste e Sul, a Mediana foi maior, igual a 57,1.
A nota Máxima 100,0 foi alcançada em todas as Grandes Regiões.
As notas mínimas ( 0,0 ) também foram iguais para todas as regiões.
Na Tabela 6.8, são apresentados o Índice de Facilidade e o Índice de Discriminação ( ponto-bisserial ) para cada uma das questões objetivas do Componente de Formação Geral.
Quanto ao Índice de Facilidade, foram usadas as seguintes cores para diferenciar o nível de dificuldade da questão :  Azul para as questões classificadas com índice Muito fácil ( \&gt ; =0,86 ), verde para as questões classificadas com índice Fácil ( 0,61 a 0,85 ), amarelo para as questões classificadas com Médio ( 0,41 a 0,60 ), vermelho para as questões classificadas com Difícil ( 0,16 a 0,40 ) e roxo para as questões classificadas com Muito difícil ( \&lt ; =0,15 ).
Já quanto ao Índice de Discriminação, foram usadas as seguintes cores para qualificar a questão :  As questões classificadas com índice Fraco receberam a cor vermelha ( \&lt ; =0,19 ), as classificadas com Médio receberam a cor amarela ( 0,20 a 0,29 ), as classificadas com Bom receberam a cor verde ( 0,30 a 0,39 ) e as classificadas com Muito bom ( \&gt ; =0,40 ) receberam a cor azul.
As questões objetivas do Componente de Formação Geral, segundo o Índice de Facilidade, foram assim avaliadas : das oito questões, nenhuma teve o Índice de Facilidade classificado como Muito fácil, e duas questão foram classificadas como Fácil, por terem índice de acertos entre 0,61 e 0,85.
Três questões foram consideradas com índice de dificuldade Médio, situando -se no intervalo entre 0,41 e 0,60 do Índice de Facilidade, ou seja, houve entre 41,0 \% e 60,0 \% de acertos, enquanto duas questões foram classificadas na categoria Difícil, situando -se no intervalo entre 0,16 e 0,40.
Por fim, uma questão apresentou menos de 15 \% de acertos, razão pela qual foi classificada como Muito difícil.
O Índice de Facilidade variou de 0,09 a 0,84, e o de Discriminação, de 0,19 a 0,53.
Na Tabela 6.9, é apresentada a distribuição das questões, considerando -se simultaneamente a classificação dos dois índices.
As seis questões com Índice de Discriminação Muito bom figuraram entre três níveis de dificuldade Fácil, Médio e Difícil : duas classificadas na categoria Fácil ( questões 4 e 6 ) do Índice de Facilidade, três na categoria Médio ( questões 2, 5 e 8 ) e, por fim, uma na categoria Difícil ( questões 7 ).
Em particular, a questão 4 foi a que apresentou o maior poder discriminatório, com índice 0,53, e foi considerada Fácil em termos de facilidade, com uma proporção de 0,61 acertos.
O máximo de acertos foi alcançado pela questão 6, com um Índice de Facilidade de 0,84.
As questões 3 e 7 foram consideradas como Difícil, com índices de facilidade 0,23 e 0,30, respectivamente.
Já os índices de discriminação foram Muito bom para a questão 7 ( 0,41 ) e Bom para a questão 3 ( 0,32 ).
A questão de número 1 apresentou um Índice de Facilidade de 0,09 ( Muito Difícil ), ou seja, menos do que 10 \% dos estudantes conseguiu resolvê -la, dentro do universo de participantes.
Seu Índice de Discriminação foi Fraco ( 0,19 ).
Por isso, a questão 1 não entrou no cômputo das notas das questões objetivas de Formação Geral dos estudantes de Licenciatura em Letras-Português e Espanhol.
No Gráfico 6.4, para exemplificar, é analisado o comportamento da questão de número 4 de Formação Geral.
Trata -se de uma questão considerada Fácil, em relação à facilidade e a que obteve o maior Índice de Discriminação dessa parte da prova.
Neste gráfico, cada uma das cinco curvas representa o percentual de respostas em determinada alternativa da questão, em função do número de acertos dos estudantes nessa parte da prova ( Formação Geral/Múltipla Escolha ), antes de possíveis eliminações pelo critério do ponto-bisserial.
A curva em verde corresponde à alternativa B, a alternativa correta para essa questão.
Observa -se que dentre os estudantes com menor número de acertos, nessa parte do exame, a situação mais frequente foi a escolha da alternativa E ( em vermelho ), incorreta.
Por exemplo, dentre os estudantes que acertaram três questões, 54,5 \% escolheram a alternativa correta B, 24,1 \% escolheram a alternativa E ( em vermelho ), 12,2 \% escolheram a alternativa D ( em roxo ), 7,5 \% escolheram a alternativa A ( azul ) e 1,7 \% a C ( preto ).
Entre os que acertaram três respostas entre as questões de múltipla escolha de Formação Geral, nenhum deixou a questão em branco ou marcou mais de uma alternativa, invalidando a questão.
À medida que o número de acertos aumenta, indicando desempenho melhor nessa parte da prova, aumenta concomitantemente a proporção de estudantes que selecionaram a alternativa correta B, atingindo 100 \% para os estudantes com sete acertos.
Essa análise permite verificar como a questão discriminou os grupos de desempenho, justificando -se o alto índice obtido na questão.
Cumpre notar que não é possível inferir deste gráfico nem o índice de facilidade ( que seria uma média da proporção ponderada pela quantidade de alunos com cada uma das notas ), nem o índice de discriminação ponto-bisserial, por razão equivalente.
No caso extremo no qual a grande concentração dos acertos dos alunos fosse abaixo de quatro, o índice de facilidade seria obrigatoriamente abaixo de 40 \% ( neste exemplo ).
Caso a concentração fosse em seis acertos ou mais, o índice seria obrigatoriamente acima de 60 \%.
Os gráficos relativos às demais questões de Formação Geral constam do Anexo I. 6.2.2 Componente de Conhecimento Específico Na Tabela 6.10, são apresentadas as Estatísticas Básicas em relação às questões objetivas do Componente de Conhecimento Específico da prova de Letras-Português e Espanhol ( Licenciatura ) por Grande Região.
Nessa parte da prova, uma questão foi anulada pela Comissão Assessora da Área, a questão 22.
Além disso, como será discutido ainda nesta seção, cinco questões objetivas do Componente de Conhecimentos Específicos não foram usadas no cômputo das notas por terem sido descartadas pelo critério do ponto-bisserial.
Assim, as notas foram calculadas com base em 21 das 27 questões objetivas de Conhecimentos Específicos.
A Média do Brasil deste componente foi 43,3.
A menor Média foi observada na região Nordeste ( 39,0 ), e a maior, na região Sul ( 48,5 ).
O Desvio padrão de todo o Brasil foi 17,9, sendo o menor Desvio padrão encontrado na região Centro-Oeste ( 15,6 ), e o maior, na região Sudeste ( 18,3 ).
A Mediana de todo o Brasil foi 42,9.
Nas regiões Norte, Nordeste e Centro-Oeste, a Mediana foi menor do que a nacional ( 38,1 ), e nas regiões Sul e Sudeste, foi maior ( 47,6 ).
A nota Máxima da prova foi 100,0, obtida nas questões objetivas do Componente de Conhecimento Específico, por, pelo menos, um aluno das regiões Nordeste e Sudeste.
Na região Norte, a nota Máxima foi 85,7, na região Sul foi 90,5 e na região Centro-Oeste foi 81,0.
A nota Mínima foi diferente de zero apenas na região Sul ( 4,8 ).
Na Tabela 6.11, são apresentados os Índices de Facilidade e Discriminação ( ponto-bisserial ) das questões objetivas do Componente de Conhecimento Específico para os estudantes de Letras-Português e Espanhol ( Licenciatura ).
Para facilitar a diferenciação das questões, serão usadas as mesmas cores da Tabela 6.8 para as diferentes classificações dos Índices de Facilidade e de Discriminação.
Observa -se que uma questão objetiva foi anulada pela CAA e as análises dizem respeito às 26 questões válidas.
A partir do Índice de Facilidade obtido, pode -se concluir que um pouco mais de metade das questões objetivas da prova foi considerado, pelo menos, Difícil : das 26 questões válidas, 14 foram classificadas como Difícil ( 12 ) ou como Muito difícil ( duas ).
Além disso, duas questões foram classificadas como Fácil, e outras dez consideradas como Médio.
Nenhuma questão foi classificada como Muito fácil.
Já quanto aos índices de discriminação das questões objetivas do Componente de Conhecimento Específico da prova, tem -se como resultado a seguinte classificação : 14 das 26 questões foram consideradas como boas, enquanto seis delas tiveram Índice de Discriminação Muito bom.
Assim, para um pouco mais de três quartos das questões – 20 em 26 – os Índices de Discriminação foram Bom ou Muito bom.
Dentre as demais, uma delas foi classificada como Médio, e outras cinco, como Fraco, sendo seis, por conseguinte, a quantidade de questões nos dois patamares mais baixos de discriminação.
Constata -se, assim, que a prova – no que se refere ao Componente de Conhecimento Específico – possuía boa capacidade de discriminar dentre aqueles que dominavam ou não o conteúdo.
Na Tabela 6.12, é apresentada a distribuição das questões considerando -se simultaneamente a classificação dos dois índices.
Dentre as seis questões que alcançaram os maiores índices de discriminação, classificadas com Índice de Discriminação Muito bom, uma delas, a de número 13, foi a de índice mais elevado ( 0,46 ), sendo classificada quanto ao Índice de Facilidade como Médio, com 58 \% dos estudantes marcando a opção correta.
A questão de número 12 foi a mais difícil dentre as 26 questões específicas válidas, com baixo Índice de Facilidade, apenas 10 \% de acertos.
Essa questão apresentou poder discriminatório muito baixo, 0,11, o que comprova ter sido a mais difícil para os estudantes.
Destaca -se, também, a questão 27, com Índice de Facilidade 0,17, considerada Difícil, já que, em termos percentuais, corresponde a 17 \% de estudantes que responderam acertadamente.
Seu Índice de Discriminação foi Fraco, igual a 0,13.
Tais questões foram, portanto, pelo critério ponto-bisserial, consideradas inadequadas.
Por isso, as questões 12 e 27, além das questões 10, 16 e 19 foram eliminadas do cômputo da nota final.
A título de exemplo das análises do comportamento das questões objetivas, no Gráfico 6.5, é analisada a questão 13 do Componente de Conhecimento Específico.
Essa questão foi considerada, pelas respostas dos estudantes avaliados na prova, como uma questão com nível de facilidade Médio, 0,58, ou seja, 58 \% dos estudantes assinalaram acertadamente a opção E, correspondente ao gabarito.
Como já comentado, seu Índice de Discriminação foi igual a 0,45, classificado como Muito bom, o maior valor de discriminação.
Neste gráfico, cada uma das cinco curvas representa o percentual de respostas em determinada alternativa da questão 13, em função do número de acertos dos estudantes nessa parte da prova, antes de possíveis eliminações de questões pelo critério do ponto-bisserial.
A alternativa correta E, representada no gráfico pela curva em vermelho, foi escolhida em maiores proporções pelos alunos com desempenho melhor nessa parte da prova.
Já as alternativas incorretas, também denominadas distratores, foram selecionadas, principalmente, por aqueles com notas mais baixas.
Observa -se que a soma não é 100 \%, por causa das questões não respondidas ou com mais de uma opção marcada.
Aqueles com nota zero, na sua quase totalidade, deixaram esta questão em branco ou marcaram mais de uma alternativa, comportamento considerado inválido.
A proporção de alunos que selecionou a resposta correta E aumenta gradativamente, chegando a atingir 100 \% para 21 acertos ou mais, enquanto a proporção dos que escolheram alternativas incorretas decai, a partir de três ou quatro acertos, em função do número de acertos nessa parte da prova.
Os gráficos relativos às demais questões do Conhecimento Específico constam do Anexo I. 6.3 ANÁLISE DAS QUESTÕES DISCURSIVAS Nesta seção, são apresentadas estatísticas selecionadas e histogramas das Questões Discursivas de Formação Geral ( 6.3.1 ) e Conhecimento Específico ( 6.3.2 ).
São também apresentadas e comparadas as médias de subpopulações, caracterizadas por Grande Região. 6.3.1 Componente de Formação Geral As análises dos resultados de desempenho dos estudantes de Letras-Português e Espanhol ( Licenciatura ) nas duas questões discursivas relativas à Formação Geral encontram -se na Tabela 6.13 e no Gráfico 6.6.
Na Tabela 6.13, observa -se que a nota Média nesse conjunto de questões foi acima da obtida nas objetivas.
Os estudantes, de todo o Brasil, obtiveram, em Formação Geral, Média 49,1 nas questões objetivas e 56,5 nas questões discursivas.
Pode -se notar também que o Desvio padrão nesse conjunto de questões foi um pouco maior do que o obtido nas objetivas, 21,0 nas questões objetivas e 23,7 nas questões discursivas.
A maior Média foi obtida na região Sul ( 61,6 ), e a menor, na região Nordeste ( 52,9 ).
A Mediana de todo o Brasil, neste componente, foi 62,0.
Nas regiões Norte, Nordeste ( 58,5, em ambas ) e Centro-Oeste ( 57,5 ), a Mediana foi menor do que a nacional e nas regiões Sudeste ( 65,0 ) e Sul ( 65,5 ) foi maior que o valor para o Brasil.
A nota Máxima ( 97,5 ) foi obtida por, pelo menos, um estudante da região Sudeste.
Nas regiões Nordeste e Sul, a nota Máxima foi 97,0, na região Norte foi 95,0, e na região Centro-Oeste foi 93,0.
A nota Mínima ( 0,0 ) foi a mesma em todas as regiões do Brasil.
No Gráfico 6.6, está representada a distribuição das notas nas questões discursivas do Componente de Formação Geral.
A moda dessa distribuição ocorre no intervalo ( 70 ; 80 ], com frequência um pouco maior do que a do intervalo ( 60 ; 70 ].
Destaca -se também o intervalo [ 0 ; 10 ] com frequência um pouco maior do que 7,5 \% do total de notas, sendo que no intervalo [ 0 ; 10 ] incluem -se, além da nota zero, a frequência de alunos que deixaram as respostas deste tipo de questão em branco.
Na sequência, os resultados verificados para cada uma das questões discursivas de Formação Geral serão apresentados, estabelecendo -se relações com os conteúdos abordados em cada uma delas.
Os comentários da Banca de docentes corretores, a respeito do observado na correção das respostas dos estudantes, suas impressões e conclusões serão apresentados junto à análise de cada questão.
Cumpre esclarecer que, tendo -se em vista que as questões discursivas de Formação Geral são padronizadas, ou seja, constam de todas as provas, os comentários da Banca são os mesmos para todas as carreiras acadêmicas, sendo direcionados a todos os estudantes que participaram do Enade/2017.
A seguir, serão analisados os desempenhos da Área de Letras-Português e Espanhol ( Licenciatura ) nas duas questões discursivas de Formação Geral do Enade/2017, comparando -se os resultados obtidos com comentários para cada questão. 6.3.1.1 Análise de Conteúdo da Questão Discursiva 1 do Componente de Formação Geral Os dados de Letras-Português e Espanhol ( Licenciatura ), obtidos a partir das respostas à questão 1, encontram -se na Tabela 6.14 e no Gráfico 6.7.
Nessa questão – de desempenho um pouco melhor entre as duas de Formação Geral – os alunos de todo o Brasil, tiveram Média 56,6.
A maior Média para a questão 1 foi obtida na região Sul ( 60,6 ), e a menor, na região Nordeste ( 53,8 ).
Quanto à variabilidade das notas, o Desvio padrão de todo o Brasil foi 27,5.
O menor Desvio padrão foi obtido na região Sul ( 24,4 ), e o maior Desvio padrão foi obtido na região Nordeste ( 28,3 ).
As medianas do Brasil, como um todo e das regiões Norte, Nordeste e Centro-Oeste foram iguais ( 60,0 ).
Nas regiões Sudeste e Sul, a Mediana foi maior, 65,0 em ambas.
As notas máximas e mínimas da questão discursiva 1 foram, sem exceção, as mesmas para todas as regiões do Brasil, respectivamente, 100,0 e 0,0.
No Gráfico 6.7, mostra -se a distribuição das notas na questão discursiva 1 do Componente de Formação Geral.
Observa -se que a maior frequência foi o intervalo ( 70 ; 80 ], a moda da distribuição.
Pouco menos de 8 \% dos alunos deixaram a resposta a essa questão em branco, um máximo local. 6.3.1.2 Comentários sobre a correção de Conteúdo das respostas à Questão Discursiva 1 O enunciado era claro e indicava um tema de extrema relevância, qual seja, a incidência de doença grave integrante do rol das DST ( Doenças Sexualmente Transmissíveis ), que encaminha para a análise do direito fundamental à saúde, que tem assento constitucional ( CRFB, art. 6º. caput : São direitos sociais a educação, a saúde, a alimentação, o trabalho, a moradia, o transporte, o lazer, a segurança, a previdência social, a proteção à maternidade e à infância, a assistência aos desamparados, na forma desta Constituição ).
A questão focalizava, por meio de três textos, um tema social extremamente relevante, qual seja, a relação entre homem e mulher nas suas interações íntimas, no caso representadas pela vida sexual, e o reflexo das relações de poder veiculadas pela sociedade, que apontam, no mais das vezes, o predomínio do homem.
Tal posição de superioridade é justificada pelas convenções arraigadas que acarretam a educação dos homens e mulheres por padrões patriarcais, que redundam em posições machistas.
Um dos textos demandava uma reflexão sobre a realidade dos casais que convivem em relações estáveis ou fugazes e não têm o controle dos efeitos das relações sexuais, que abrangem a gravidez não planejada e, em decorrência do não cuidado, a inserção de doenças transmitidas sexualmente.
Esse \textbf{panorama} é refletido na pouca influência das mulheres sobre os homens quanto à proteção, normalmente representada pelo uso de preservativos nas relações intimas, e após a constatação da doença, no encaminhamento das pessoas para tratamento.
Como efeito colateral, está a transmissão de doenças para os fetos.
Outro texto trazia à tona o maior cuidado que as mulheres têm, em geral, com sua saúde, notadamente no momento em que se descobre grávida e que redunda na realização de exames obrigatórios de pré-natal, que podem facilmente detectar doenças e, caso encontradas, ter o seu tratamento encaminhado.
A par disso, destaca -se que a crise de recursos financeiros agravou a situação, que já recebe influxos sociais negativos, como a dificuldade do homem em utilizar preservativos, realizar exames e procurar tratamentos quando adoece.
Como muito divulgado pelas mídias, a crise financeira do país repercutiu diretamente na prestação de serviços de saúde em todas as esferas : distrital, federal, estadual e municipal.
Assim, agravou -se o problema da falta de medicamentos essenciais para o tratamento da sífilis e, certamente, de outras doenças.
Como se verifica, a temática indicada no enunciado, representada pelos textos, pertence à realidade conhecida amplamente pelos estudantes, pois também é divulgada por todos os meios de comunicação.
Trata -se de tema geral que, no entanto, normalmente, não se inclui nas grades curriculares de cursos superiores que não sejam da área de Saúde.
Apesar disso, a temática é de conhecimento geral e relacionada aos direitos humanos, esses sim, ponto de integração com as diretrizes curriculares.
A linguagem utilizada era clara, indicando a existência dos fatos vinculados a questões relevantes e à saúde pública.
No entanto, como os textos trazem uma ampla gama de possibilidades de interpretação, ocorreram algumas compreensões diversas das esperadas.
Além disso, ressalta -se que a amplitude dos textos disponibilizados no enunciado fornece elementos que poderiam ser simplesmente reproduzidos pelos concluintes nas suas respostas.
O enunciado da questão encaminha a resposta para dois aspectos.
O estudante deveria abordar conteúdo pertinente à vulnerabilidade das mulheres às DST e o papel social do homem em relação à prevenção dessas doenças.
Para responder a esses aspectos, um número expressivo de concluintes utilizou somente os elementos dos textos motivadores, sem nada acrescentar.
A resposta deveria conter também uma segunda parte, a indicação de duas ações voltadas para o público masculino, a serem adotadas no âmbito das políticas públicas de saúde ou de educação, para a redução do problema.
Verificou -se que o desempenho dos estudantes nessa questão foi de médio para bom.
Ao se analisarem as respostas, constata -se que muitos foram os que apresentaram os tópicos presentes no padrão de resposta, amplo o suficiente para albergar as respostas apresentadas pelos concluintes.
Tal resultado pode ser justificado, tendo em vista tratar -se de tema bastante difundido nas diversas mídias, por ter sido apresentado de forma ampla pelos textos motivadores presentes no enunciado.
Por outro lado, o debate no ambiente universitário sobre as relações de poder entre homens e mulheres tende a ser marcante pelo movimento de empoderamento feminino, liderado por organizações sociais com ampla atuação nos estabelecimentos de nível superior.
De forma recorrente, as relações patriarcais entre homens e mulheres foram vinculadas ao modelo machista vigente.
Essas condições, segundo os concluintes, explicam a incidência da doença e a situação de vulnerabilidade das mulheres, bem como a função social dos homens.
Outras respostas apontaram o poder público, nas áreas de saúde e/ou educação, como órgãos responsáveis pela solução do problema.
Para os que se fixaram na área da saúde, foram frequentes a indicação da necessidade de mais medicamentos, vacinas, educação sexual, médicos e postos de atendimentos.
Alguns, no entanto, apenas indicaram ações individuais para o homem, e outros, ações individuais para as mulheres.
Também figuraram respostas que apresentavam ações da igreja como solução, além de outras que faziam referências a campanhas no seio das famílias. 6.3.1.3 Análise de Conteúdo da Questão Discursiva 2 do Componente de Formação Geral Na Tabela 6.15, percebe -se que o desempenho médio dos estudantes na questão discursiva 2 ( média 55,3 ) foi um pouco inferior ao obtido na questão discursiva 1 ( média 56,6 ).
A região Sul foi aquela cuja Média, também nessa questão, foi maior ( 61,7 ), e a de menor Média foi a região Nordeste ( 50,8 ).
Quanto à variabilidade das notas, o Desvio padrão de todo o Brasil foi 32,0, superior ao obtido na questão discursiva 1 ( 27,5 ).
O maior desvio nessa questão foi obtido na região Nordeste ( 33,0 ), enquanto o menor foi obtido na região Sul ( 29,6 ).
A Mediana de todo o Brasil foi 70,0, a mesma das regiões Sudeste e Sul.
Nas regiões Norte, Nordeste e Centro-Oeste, as medianas foram iguais a 50,0.
As notas máximas ( 100,0 ) e mínimas ( 0,0 ) também foram as mesmas em todas as regiões, sem exceção.
No Gráfico 6.8, apresenta -se a distribuição das notas na questão discursiva 2, do Componente de Formação Geral.
Observa -se que a maior frequência corresponde à dos alunos que obtiveram nota dentro do intervalo ( 70 ; 80 ].
Destacam -se, também, os alunos que deixaram a resposta dessa questão em branco e os alunos que obtiveram nota dentro dos intervalos ( 20 ; 30 ], ( 50 ; 60 ] ou ( 90 ; 100 ], caracterizando -se como máximos locais.
Nota -se, ainda, que as notas ficaram mais espalhadas em comparação às notas da questão discursiva de número 1, o que pode ser constatado também pela comparação do desvio padrão das notas da questão discursiva 2 ( 32,0 ) e da questão discursiva 1 ( 27,5 ). 6.3.1.4 Comentários sobre a correção de Conteúdo das respostas à Questão Discursiva 2 O enunciado era claro e abordava um tema de extrema relevância : a questão de gênero.
Trata -se de uma proposta polêmica, pois abrange a análise da inserção social de minoria vinculada à questão de gênero.
O enunciado solicitava que o concluinte discorresse sobre a importância do nome para as pessoas transgêneras e propusesse uma medida, no âmbito das políticas públicas, para facilitar o acesso dessas pessoas à cidadania.
A respeito da primeira parte, a importância do nome, esperava -se que o estudante abordasse as dificuldades que o uso do nome civil ( nome de registro, nome de batismo, nome da certidão de nascimento, por exemplo ) pode acarretar para as pessoas transgêneras, além de demonstrar conhecimento sobre a importância do nome social, direito já reconhecido através, principalmente, de decisões judiciais e atos administrativos, como decretos e portarias que dispõem sobre o uso do nome social.
Em atendimento à segunda parte do comando, o participante deveria propor uma política pública associada ao direito à cidadania das pessoas transgêneras.
Sobre esse aspecto, esperava -se a compreensão do significado de política pública, ou seja, propostas de âmbito governamental, e não, iniciativas de Organizações Sociais, igrejas ou grupos sociais de qualquer tipo.
O nascimento das pessoas, no Brasil, é objeto de inscrição em cartórios de Registro Civil das Pessoas Naturais, órgão responsável por estabelecer a certificação pública dos dados de todos os indivíduos que não somente nascem, mas casam, eventualmente se divorciam e vêm a óbito.
Daí o nome civil.
Para flexibilizar a correção, outras nomenclaturas foram consideradas, a saber : identificação de nascimento, nome de batismo, nome de registro, nome próprio, nome verdadeiro, que diferenciam o nome civil daquele que o indivíduo deseja utilizar para traduzir a sua opção pessoal de gênero, que se convencionou denominar nome social.
Quanto ao entendimento do termo transgênero, deve -se destacar que o movimento de mudança do gênero é contemporâneo e não está necessariamente vinculado à mudança de sexo.
Existem várias hipóteses em jogo.
Uma é a que envolve cirurgias de mudança de sexo ( a menina extirpa os seus seios e inclui órgão genital por intervenção cirúrgica ; o menino inclui seios e extirpa o seu órgão genital e inclui o órgão feminino também por intervenção cirúrgica ).
Há outro movimento que envolve apenas as vestimentas e o linguajar corporal, além dos relacionamentos ( as meninas assumem o trajar, o andar, os gostos dos meninos e mantêm relacionamentos afetivos com meninas ; os meninos assumem o trajar, o andar, os gostos das meninas e mantêm relacionamentos afetivos com meninos ).
Em relação à discussão da primeira solicitação do comando da questão, destaca -se que todos têm nome civil, que foi escolhido pelos seus pais ou responsáveis e, na falta deles, por autorização do Juiz de Direito, pautado no sexo detectado no nascimento.
Já o nome social é aquele que indica a preferência do indivíduo na sua identificação nas relações sociais e que destoa do seu gênero de registro civil.
Ressalte -se que o Brasil não possui ainda legislação específica sobre pessoas transgêneras.
Observa -se um relativo progresso no reconhecimento de direitos, através, principalmente, de decisões judiciais e atos administrativos, como decretos e portarias que dispõem sobre o uso do nome social.
Sem uma lei que defina os procedimentos da alteração dos documentos para pessoas transexuais, essa parcela da população LGBT ( lésbicas, gays, bissexuais, transexuais e transgêneros ) é obrigada a procurar na Justiça o reconhecimento de sua identidade, em processos que podem ser longos e que dependem do olhar de determinados, profissionais, como médicos, psicólogos, assistentes sociais, bem como do entendimento dos advogados, defensores públicos, juízes e do Judiciário, como um todo.
Atualmente, tramita na Comissão de Direitos Humanos da Câmara dos Deputados o Projeto de Lei João Nery ( 5002/2013 ), dos deputados Jean Willys ( PSOL-RJ ) e Erika Kokay ( PT-DF ), que determina que o reconhecimento da identidade de gênero seja um direito do cidadão.
O projeto recebeu o nome do primeiro trans-homem operado no Brasil.
Grande parte dos respondentes considerou a incongruência entre a imagem corporal, representada pelo gênero assumido pelos trans, e o nome que apresenta nos seus documentos como causa de sérios constrangimentos e embaraços e as dificuldades para que esses sujeitos acessem direitos básicos, como saúde, educação e empregabilidade em uma sociedade heteronormativa e preconceituosa.
Muitos demonstraram reconhecer que a identidade e expressão de gênero, através do nome social, pressupõe que sejam respeitadas a dignidade da pessoa humana, a intimidade, a liberdade e o direito à cidadania plena.
No entanto, também foram encontradas, em quantidades significativas, respostas que demonstravam a existência de uma interpretação social conservadora que não admite a mudança de sexo nem do nome social.
Existe um acentuado preconceito contra todos os movimentos que buscam identificação social diversa da civil.
Isso inclui os variados grupos : gays, travestis, transgêneros, dentre outros.
Já em relação ao segundo tópico a ser abordado pelo concluinte, é notório que as pessoas trans encontram -se destituídas de direitos no campo das políticas públicas, como saúde, habitação, emprego, educação, justiça, dentre outros.
Nesse sentido, uma grande diversidade de medidas, no âmbito das políticas públicas, poderia ser proposta.
Ações voltadas para a capacitação de profissionais da educação, saúde, segurança pública, dentre outras áreas, para melhor atender às demandas e especificidades desse público alvo seriam necessárias.
Ainda em relação à educação, seriam também necessárias ações que contribuíssem para diminuir a evasão, pois o ambiente escolar é difícil para uma pessoa transgênera, por conta das situações discriminatórias e preconceituosas vindas dos colegas de classe, e mesmo das/dos professores e/ou de outros profissionais da área, não podem ser mais um complicador nesse processo.
Ainda no que tange ao assunto, percebe -se que, ao serem excluídas dos bancos escolares, as oportunidades no mercado de trabalho formal praticamente inexistem.
Um outro aspecto que se evidencia no cotidiano desse segmento da sociedade diz respeito às situações de violência sofridas pelas pessoas transgêneras.
Nesse sentido, algumas propostas foram elencadas para minimizar o problema e combater a violência.
Políticas envolvendo melhor formação de funcionários das polícias civil e militar, por exemplo, somariam esforços na qualificação das abordagens policiais e dos atendimentos prestados pelas forças de segurança, pois o respeito à orientação sexual e à identidade e expressão de gênero é, antes de tudo, o respeito à dignidade, à cidadania e à própria democracia.
Nesse sentido, revela -se essencial que, ao sofrerem agressões, sejam elas físicas, psicológicas ou morais no espaço público ou domiciliar, as pessoas trans possam ser atendidas, com respeito inclusive ao nome social, nas delegacias comuns e, inclusive, nas Delegacias Especiais de Atendimento à Mulher ( DEAM ), pois esse é o órgão institucional que possivelmente detém as qualificações necessárias para o atendimento ao gênero feminino ao qual as travestis e mulheres trans efetivamente pertencem.
Também ocorre uma latente inacessibilidade das pessoas trans à assistência jurídica.
Nesse sentido, é necessário implementar projetos de capacitação para profissionais do poder judiciário, como um todo, de modo a prestarem a assistência jurídica gratuita na defesa intransigente dos direitos sociais, civis, trabalhistas e políticos dessa população.
Dessa maneira, a Justiça agiria de forma mais acessível e sensível para o segmento, que terá suas especificidades observadas a começar pelo devido respeito do direito ao uso do nome social.
Nesse âmbito, o mais comum foi propor a adoção de medidas punitivas para quem viola o direito à autodeterminação de gênero e a elaboração de leis que garantam a mudança do nome e assegurem outros direitos para as pessoas trans.
Propostas envolvendo a formação de profissionais, como as mencionadas aqui, e de conscientização da população foram as mais frequentes.
Muitas outras enfatizavam o direito ao uso respeitado do nome social e a garantia de alteração do registro civil, independentemente da cirurgia de redesignação sexual ou transgenitalização.
Muitos questionaram o fato de não bastar a autodeclaração no processo de alteração de nome, enfatizando que o nome social é uma questão de princípio de dignidade e cidadania.
Apesar de pouco mencionada nas respostas corrigidas, políticas envolvendo a promoção da saúde são muito necessárias.
Cabe aos gestores de saúde ações direcionadas, não somente à prevenção e tratamento do HIV/Aids, mas também políticas de saúde que contemplem as especificidades dessa população, tais como o Ministério da Saúde, especialmente o SUS ( Sistema Único de Saúde ), incluir, de fato, o atendimento para pessoas trans para acompanhamento da terapia hormonal, que muitas demandam, respeitando -se o direito de utilização do nome social.
De modo geral, são muitas as políticas públicas a que precisam ser implantadas.
Urge a necessidade de uma agenda mais compromissada por parte dos gestores do nosso país com essa comunidade, que ainda é muito discriminada.
Logo, tais políticas devem contemplar suas necessidades básicas como o direito de acesso aos estudos, à capacitação para o mercado de trabalho, à profissionalização, ao acesso a bens e serviços de qualidade em saúde, habitação, segurança, cultura e assistência social e, certamente, o respeito ao uso do nome social, o que também contribui para essa inclusão.
Assim, as opções apresentadas no padrão de resposta funcionaram como exemplos de políticas públicas possíveis, já que seria difícil apresentar uma lista suficientemente ampla para abranger a grande variedade de respostas encontradas.
O desempenho dos estudantes nessa questão foi médio.
As respostas, de modo geral, indicaram a necessidade de uma atualização por parte dos concluintes quanto a temas essenciais e, no caso da presente questão, a análise sobre aspectos fundamentais da questão de gênero.
Muitos, possivelmente motivados pelos textos apresentados no enunciado, restringiram -se a discorrer sobre preconceito de gênero ou mesmo sobre o preconceito de modo mais geral, sem responder ao que foi solicitado no comando da questão.
Outros apresentaram discursos religiosos ou ideológicos, também sem tratarem do que foi solicitado. 6.3.1.5 Análise de Língua Portuguesa das Questões Discursivas do Componente de Formação Geral Os dados de Letras-Português e Espanhol ( Licenciatura ), obtidos a partir das respostas às questões discursivas do Componente de Formação Geral, no que tange à Língua Portuguesa, encontram -se na Tabela 6.16 e no Gráfico 6.9.
Nesse aspecto, os alunos, de todo o Brasil, obtiveram Média 59,6.
A maior Média em relação à Língua Portuguesa foi obtida na região Sul ( 63,9 ), e a menor, na região Nordeste ( 56,3 ).
Quanto à variabilidade das notas, o Desvio padrão de todo o Brasil foi 20,7.
O menor Desvio padrão foi obtido na região Sul ( 17,7 ), e o maior Desvio padrão foi obtido na região Nordeste ( 21,6 ).
A Mediana das notas de Língua Portuguesa foi 65,0 para o Brasil, como um todo.
Nas regiões Norte, Nordeste e Centro-Oeste, a mediana foi 62,5, menor do que a nacional, e nas regiões Sudeste e Sul, foi 67,5, maior que a do Brasil como um todo.
A nota Máxima para todo o Brasil foi 95,0, com, pelo menos, um aluno tirando essa nota na região Nordeste.
Na região Norte a nota Máxima foi 85,0, e nas regiões Sudeste, Sul e Centro-Oeste, foi 92,5.
Já a nota Mínima foi zero em todas as regiões do país.
No Gráfico 6.9, mostra -se a distribuição das notas de Língua Portuguesa do Componente de Formação Geral.
Observa -se que a maior frequência ( mais do que 30 \% ) corresponde à dos alunos que obtiveram nota no intervalo ( 60 ; 70 ].
Observa -se, também, a classe dos alunos que deixaram as respostas de ambas as questões em branco representando um pouco mais do que 5 \% do total. 6.3.1.6 Comentários sobre a correção das respostas de Formação Geral em relação à Língua Portuguesa Os enunciados apresentados em 2017 não explicitaram a exigência de elaboração de um texto “ dissertativo ”, como nos anos anteriores.
O participante deveria fazer uma exposição de seus conhecimentos sobre cada um dos assuntos e estruturar seus textos de acordo com as características do registro formal adequado à situação comunicativa – avaliação de conhecimentos.
Essa configuração determina exigências quanto aos seguintes aspectos : adequação da seleção vocabular, desenvolvimento do conteúdo, organização lógica das ideias, estruturação sintática dos períodos, utilização de procedimentos de encadeamento textual e de referenciação, obediência às exigências morfossintáticas próprias da modalidade escrita da norma-padrão, respeito às regras ortográficas e às regras de acentuação gráfica.
O padrão de resposta utilizado na avaliação das questões 1 e 2 considerou os aspectos relevantes ao bom desempenho linguístico como competências distintas, de modo a permitir um mapeamento detalhado do domínio dos recursos disponíveis na Língua Portuguesa para a comunicação escrita formal.
Com base nesse objetivo, foram avaliados os seguintes aspectos : ( a ) estruturação textual condizente com o gênero solicitado e o modo de organização textual expositivo adequado ao gênero – essa competência envolve :  estruturação sintática condizente com o padrão da modalidade escrita formal da língua portuguesa de modo a garantir a clareza necessária ;  distribuição do conteúdo do texto em parágrafos, de modo a garantir a sua organização temática ;  utilização de operadores discursivos que contribuam para a progressão temática do texto, estabelecendo relações lógicas entre as ideias apresentadas, tanto do ponto de vista intrafrasal, como do interfrasal ;  utilização de procedimentos de referenciação lexical e pronominal que permitam a retomada de referentes textuais ;  utilização de sinais de pontuação que contribuam para a organização lógica da frase e do texto ;  inteligibilidade relacionada ao atendimento das exigências de estruturação textual.
Espera -se, portanto, que o participante recorra a procedimentos linguístico-discursivos para organizar seu texto, permitindo o encadeamento lógico entre suas partes de forma a garantir a progressão e a coerência textuais.
Isso significa que os seguintes procedimentos foram considerados inadequados, de acordo com o padrão de resposta proposto :  elaboração de frases fragmentadas que comprometam a estrutura lógico-gramatical do texto ;  sequência justaposta de ideias sem encaixamentos sintáticos, reproduzindo hábitos da oralidade ;  elaboração de frase com apenas oração subordinada, sem oração principal ;  emprego equivocado do conector ( preposição, conjunção, pronome relativo, alguns advérbios e locuções adverbiais ) comprometendo a expressão da relação lógica entre duas ideias, com prejuízo da clareza do texto ;  emprego do pronome relativo sem a preposição, quando obrigatória ;  repetição ou substituição inadequada de palavras sem o emprego dos recursos oferecidos pela língua ( pronome, advérbio, artigo, sinônimo ) ;  emprego inadequado dos pronomes relativos “ cujo(a) ” e “ onde ” ;  utilização inadequada dos sinais de pontuação, comprometendo a clareza textual ;  inteligibilidade relacionada ao atendimento das exigências de estruturação textual. ( b ) respeito às convenções ortográficas da norma-padrão da Língua Portuguesa – essa competência envolve o domínio das regras de acentuação gráfica e da grafia padrão das palavras ( com ausência de abreviaturas próprias da linguagem da internet ), de acordo com as convenções estabelecidas pela legislação em vigor e consubstanciadas no Vocabulário Ortográfico da Língua Portuguesa, editado pela Academia Brasileira de Letras. ( Vale notar que, nessa edição de 2017, não foi aceita a legislação anterior, no caso das regras relativas ao uso do hífen e da acentuação gráfica ).
Espera -se que o participante :  grafe corretamente as palavras ;  respeite as regras de acentuação gráfica ;  empregue maiúsculas em início de frase, em nomes próprios de pessoas, lugares ou instituições ;  grafe as siglas com letras maiúsculas ou apenas com a primeira letra maiúscula quando formarem uma palavra ;  evite abreviações como p/, vc, tb, pq, tá, né, usadas muitas vezes em escrita informal e na internet ;  obedeça às regras de separação de sílabas no final da linha. ( c ) domínio dos diferentes aspectos morfossintáticos próprios da modalidade escrita formal da norma-padrão da Língua Portuguesa – essa competência envolve : a concordância nominal, a concordância verbal, a regência nominal, a regência verbal, a flexão nominal, a flexão verbal, a correlação entre os tempos verbais, a colocação pronominal e a utilização de sinais de pontuação que contribuam para a organização lógica da frase e do texto.
Espera -se que o participante :  flexione o verbo para estabelecer concordância de número com o sujeito da frase ;  flexione o artigo, o adjetivo e o pronome para concordar em número e em gênero com o substantivo a que se referem ;  observe a regência nominal e a verbal, utilizando a preposição adequada depois de um substantivo, um verbo ou um adjetivo ;  empregue adequadamente o acento grave indicador da crase entre uma preposição e um artigo ( a+a ) ;  obedeça às regras de colocação pronominal ( próclise e ênclise ), distintas dos hábitos da oralidade ou da escrita informal ( exigência de próclise com termo atrator, não exigência de que o pronome oblíquo se ligue ao verbo auxiliar por meio de hífen ) ;  flexione adequadamente verbos, substantivos, adjetivos e pronomes no que diz respeito à expressão das categorias gramaticais ;  flexione os verbos para expressar a correlação de modo e tempo nas estruturas subordinadas.
Com base nesses critérios, foram considerados como desvios de caráter morfossintático, e não como desvios ortográficos, as alterações que envolvem mudança de classe gramatical ou de forma flexional do verbo :  eliminação da marca de infinitivo ( -r- ) e substituição por acento agudo ou ausência total de marca do infinitivo ;  confusão entre “ ão ” e “ am ” nas formas verbais ;  confusão entre “ há ” e “ a ” ;  uso de hífen para separar pronome átono – tanto uso indevido quanto omissão ( exemplo : “ esperasse ”, em lugar de “ espera -se ” ; “ falar -mos ”, no lugar de “ falarmos ” ;  “ esta ” ( no lugar de “ está ” ) ; “ mais ” ( no lugar de “ mas ” ) ; “ e ” ( no lugar de “ é ” ).  verbos “ ter ” e “ vir ” que, na terceira pessoa do plural, não apresentarem o acento circunflexo, serão considerados como desvio de concordância, em aspectos morfossintáticos. ( d ) seleção vocabular adequada à modalidade escrita formal da Língua Portuguesa, exigida pela situação comunicativa – essa competência envolve a precisão na seleção/utilização do vocabulário relacionado à temática solicitada pela questão ; a ausência de marcas da oralidade, como termos de sentido muito genérico ( “ coisa ”, “ negócio ”, “ você ” ) e termos de registros mais informais ( como gírias, jargões, frases feitas, ditados populares, termos regionais ).
Assim, espera -se que o participante respeite a adequação vocabular não usando gírias ou expressões coloquiais, evite repetição desnecessária de palavras e utilize um vocabulário mais formal, como solicitado por um texto dissertativo.
Observações : - A inteligibilidade relacionada ao atendimento das exigências de estruturação textual foi avaliada na segunda competência ( relativa aos aspectos textuais ). - Os problemas de coerência textual provocados por uso indevido do vocabulário foram avaliados na terceira competência ( relativa aos aspectos morfossintáticos e vocabulares ). - Não foram considerados como desvios morfossintáticos os problemas de caligrafia ( -a/ -o, -s, -r ). - Não foram considerados os textos grafados integralmente em caixa alta, já que não seria possível distinguir alguma marcação especial para as letras em início de frase. - Cada desvio foi considerado como uma ocorrência, mesmo que dois desvios fossem relativos ao mesmo aspecto linguístico.
A grade de avaliação do desempenho linguístico considerou, portanto, três grandes grupos de competências, segundo os aspectos explicitados anteriormente : a ) domínio das convenções ortográficas : grafia de vogais e consoantes, uso de maiúsculas e minúsculas, emprego do hífen e acentuação gráfica ; b ) domínio dos procedimentos de estruturação textual do ponto de vista microestrutural : organização interna dos períodos, emprego de conectores para a articulação lógica entre os períodos e entre os parágrafos, emprego de marcas de referenciação lexical e pronominal ; utilização dos sinais de pontuação que contribuem para a organização lógica da frase ; c ) domínio das regras de caráter morfossintático estabelecidas como modelares do ponto de vista da modalidade escrita formal da norma-padrão da Língua Portuguesa : concordância nominal e verbal, regência nominal e verbal, colocação pronominal, flexão nominal e verbal, correlação entre tempos e modos verbais, ausência de marcas de oralidade.
A seleção vocabular adequada à modalidade escrita formal da Língua Portuguesa foi incorporada a essa última competência, tendo em vista a intersecção entre as duas do ponto de vista das exigências do registro formal da modalidade escrita da norma-padrão.
Apreciação geral do desempenho dos estudantes : A correção revelou desempenhos distintos dos participantes, marcados pela falta de repertório cultural da maioria dos participantes : - a questão 1 propiciou a oportunidade de maior desenvolvimento do tema solicitado, gerando consequências na estruturação textual, já que os textos foram mais longos, mais elaborados e fluentes.
Entretanto, em virtude da diversidade de abordagens dos textos motivadores ( a falta de penicilina, a recusa no uso da camisinha pelos homens, a maior frequência das mulheres na realização de exames ginecológicos, a transmissão das DST aos bebês durante a gestação ), observaram -se respostas que se configuraram como verdadeiras paráfrases, com comprometimento do caráter autoral desejável em uma questão discursiva de Formação Geral.
Alguns participantes selecionaram um trecho de cada texto motivador para compor suas respostas, outros copiaram trechos inteiros. - a questão 2 teve um comportamento distinto, devido ao caráter polêmico que o caracteriza.
Assim, os textos são mais curtos, com vocabulário repetitivo e preso aos textos motivadores.
Além disso, observou -se grande ocorrência de protesto em função do tema, ora por motivos religiosos ora por preconceito de gênero.
Quanto aos aspectos linguísticos analisados durante esta avaliação, os resultados observados estão descritos nos parágrafos seguintes.
Aspectos ortográficos : O desempenho dos participantes revelou uma diferença muito grande nos dois aspectos analisados nesta competência : baixo índice de desvios da grafia padrão e grande índice de desvios de acentuação.
Em vários casos, ocorre ausência completa de acentuação gráfica.
Os resultados revelam que a tendência dominante entre os universitários brasileiros é a eliminação da acentuação gráfica, talvez motivada pelos hábitos relacionados às redes sociais e pela ausência de esclarecimento dos meios de comunicação, das autoridades e das escolas sobre as decisões do Acordo Ortográfico da Língua Portuguesa de 1990.
Os casos mais sistemáticos de eliminação do acento indicador da sílaba tônica são :  palavras proparoxítonas ( “ sifilis ”, “ proximos ”, “ politicas ”, “ publicas ”, “ transgeneros ” ) ;  palavras paroxítonas terminadas em ditongo crescente ( “ necessario ”, “ noticiarios ”, “ individuo ”, “ dependencia ”, “ varios ”, “ propria ”, “ transmissiveis ”, “ ocorrencia ” ) ;  palavras paroxítonas com hiato ( “ saude ” )  palavras oxítonas ( “ ninguem ”, “ esta ”, “ ate ”, “ tambem ”, “ prevencao ” ).
Por outro lado, destaca -se o uso indevido do acento gráfico em palavras como “ genêro ”, “ melâncolia ”, “ prevênção ”, “ intervênção ”, por exemplo.
Há, também, uma tendência a acentuar paroxítonos como se fossem oxítonos : “ transmissivéis ”.
Quanto ao domínio das convenções relativas à grafia das palavras, observam -se desvios recorrentes, como a hipercorreção pela escolha de “ e ” no lugar de “ i ”, por influência de hábitos da oralidade ( “ descriminação ” no lugar de “ discriminação ”, “ entervenção ” no lugar de “ intervenção ” ) ; desvios de grafia relacionados à variação diastrática podem ser observados em “ estrupo ”, “ subjulgadas ”, “ precoseito ”, “ soubre ”, “ vecendo ”, “ indesencia ”, “ apolojia ”, “ fulga ”, “ dereitos ”, “ sifelis ”.
Observam -se, também, casos de inadequação no uso da maiúscula para destacar determinadas palavras-chave do texto, como “ Transgêneros ”, “ Brasileiros ”, “ Homens ”.
Destaque -se, também, o grande número de participantes que grafam os textos inteiramente em caixa alta.
Vale observar, também, que, ao contrário do que se esperava, não apareceram abreviaturas próprias do “ internetês ”, relacionadas ao uso de redes sociais e e-mails.
Aspectos textuais : Esta competência é a que se revela como a mais problemática entre os participantes, porque são muitos os problemas observados, desvios acumulados durante toda a formação do estudante e que não se resolvem com um estudo autodidata, como acontece com regras ortográficas ou morfossintáticas.
São eles : sequência justaposta de ideias sem encaixamentos sintáticos ; redução drástica de estruturas subordinadas, ao lado do aumento na frequência de estruturas coordenadas e absolutas ; redução no uso de conectores para expressar relações lógicas essenciais à construção do texto, substituídas pela exigência de inferência por parte do interlocutor para suprir a sua ausência ; emprego equivocado de operadores que não estabelecem relações lógicas coerentes entre ideias do texto ; emprego inadequado do pronome relativo ( com omissão da preposição ou a utilização de pronome inadequado, como “ onde ” ) ; repetição exaustiva de termos sem a utilização de procedimentos mais sofisticados de substituição ( hiperonímias, hiponímias, nominalizações, expressões metafóricas ) ; frases fragmentadas que comprometem a estrutura lógico-gramatical ; frases formadas apenas por oração subordinada, sem oração principal.
Um importante aspecto a destacar é o baixo desempenho de uma parte dos participantes em relação à estrutura formal do texto produzido, o que é preocupante ao se levar em conta que são graduandos em fase final de formação.
São frequentes os casos de desvios de estruturação frasal, com uso inadequado ou ausência de conectivos entre parágrafos e entre frases.
Em uma parte dos textos, falta textualidade e domínio do registro padrão da língua.
Na verdade, observam -se relações linguísticas quase agramaticais, como as estabelecidas pela sequência de gerúndios sem o apoio de um ponto de partida para a organização das informações gramaticais e semânticas.
Observou -se que uma parte dos participantes não distribuiu as ideias em parágrafos, talvez devido ao pequeno número de linhas disponibilizadas para a resposta da questão ou, quem sabe, pela suposição de que não seria necessária essa divisão por não se tratar de um texto no modelo de uma redação dissertativo-argumentativa, como solicitado nos vestibulares.
Em função do encaminhamento dos enunciados das duas questões, que solicitaram uma análise do problema e encaminhamentos de políticas públicas, houve uma grande tendência, também, de construção de dois parágrafos desconexos, sem utilização de elementos coesivos adequados à progressão textual.
Quanto à utilização dos mecanismos de referenciação, deve -se destacar a ocorrência de repetições de palavras ou expressões sem a utilização de termos sinônimos ou pronomes, como seria adequado.
Quanto à utilização dos sinais de pontuação, observou -se uma grande precariedade nos textos analisados.
É muito frequente a ocorrência de parágrafos sem marca interna de pontuação para separar os períodos.
Vale observar que não foi penalizada a ausência de vírgula para destacar locuções ou adjuntos adverbiais de pequena extensão deslocados de posição na frase, por ser um uso opcional.
São os seguintes os tipos de problemas encontrados : a ) vírgula : utilização de vírgula para separar o sujeito e o predicado ; ocorrência de apenas uma das vírgulas para separar uma palavra, uma expressão ou uma oração encaixada ; uso de vírgula no lugar do ponto para separar ideias que constituem períodos distintos ; ausência de vírgula para separar elementos de uma enumeração ; ausência de vírgula para separar oração adjetiva explicativa ou utilização inadequada para separar oração adjetiva restritiva ; b ) ponto e vírgula : utilização do ponto e vírgula no lugar de vírgula ; c ) ponto final : ausência de ponto final para separar períodos.
Aspectos morfossintáticos e vocabulares : Em relação à regência, o desvio mais frequente é a falta do sinal indicativo da crase – isso revela que o usuário não tem consciência de que, sob a forma do termo “ a ”, existe a presença de uma contração entre a preposição “ a ” ( exigida pela regência do termo anterior ) e o artigo definido “ a ”.
Um desvio de regência significativo, nos últimos anos, é a utilização inadequada de uma preposição ou sua ausência após o verbo ou o nome ( substantivo ou adjetivo ).
Outro problema relacionado à regência verbal e à nominal, encontrado frequentemente nas questões, foi a ausência de preposição antes de pronome relativo, processo generalizado na modalidade oral da língua, em situações de registro informal.
Apesar da possibilidade de que essa alteração de regência se generalize no padrão escrito da Língua Portuguesa, como já está ocorrendo até em textos jornalísticos, o não emprego da preposição foi considerado inadequado neste processo de avaliação.
A concordância verbal e a concordância nominal apresentam alguns desvios muito frequentes.
Quanto à concordância de número, observou -se ausência de marca ( com sujeito anteposto ou posposto ) ou uso indevido ( uso inadequado da marca de plural comandado pelo núcleo plural da locução adjetiva, apesar de o substantivo que funciona como núcleo do sintagma nominal estar no singular ).
Uma ocorrência que se destacou foi a ausência de acento circunflexo na forma plural do presente do indicativo do verbo “ ter ”, que foi considerada como um desvio na concordância verbal e não na acentuação gráfica.
Quanto à concordância de gênero, vários casos foram observados, normalmente no âmbito de sintagmas nominais longos, em que o adjetivo está afastado do substantivo.
Deve -se destacar o aparecimento da marca de plural em verbos ou adjetivos comandados por adjunto adnominal plural, apesar de serem relacionados a núcleos substantivos no singular, evidenciando um processo de hipercorreção ( exemplo : “ O nome social dos trangêneros podem ser estabelecidos por uma legislação específica ” ).
Quanto à questão da colocação pronominal, foram poucos os desvios observados.
Concluiu -se que, no registro escrito formal, a maioria dos participantes já incorporou regras como a não introdução da frase por um pronome oblíquo e a próclise na presença de um termo atrator.
Não se adotou, entretanto, o padrão excessivamente formal descrito pelas gramáticas normativas em relação à posição do pronome oblíquo em locuções verbais, já que esse uso está muito distante da prática cotidiana, até em textos mais formais.
Quanto aos aspectos vocabulares, alguns tipos de inadequação foram observados : expressões da oralidade ; seleção vocabular incompatível com o contexto, gerando falta de inteligibilidade ; falta de domínio de vocabulário mais abstrato e de maior complexidade, essencial ao desenvolvimento do texto de base dissertativa.
O principal aspecto observado foi a excessiva repetição de certas palavras, revelando limitação de repertório vocabular.
O termo “ pessoa ”, por exemplo, chega a ser repetido até 7 ou 8 vezes em um mesmo texto.
Várias marcas de oralidade foram identificadas, embora não com alta frequência : uso do pronome relativo “ onde ” como relativo universal, falta de artigo definido antes de substantivo, falta de artigo definido antes de substantivo, reduções como “ tá ”, “ pra ”, “ pro ”, “ prum ”, expressões informais, eliminação de preposições.
Em função do tema solicitado na questão 1, a sigla DST foi grafada de diferentes maneiras : DST´s, DSTs ou DSTS.
Além disso, houve flutuação de gênero gramatical nos sintagmas que envolviam a sigla : “ os diferentes DSTs ” ou “ as diferentes DSTs ”.
Em função do tema solicitado na questão 2, houve flutuação de gênero gramatical nos sintagmas que envolviam a sua palavra-chave : “ os trans ” ou “ as trans ” ; “ os transgêneros ” ou “ as transgêneras ”.
Todas essas formas foram aceitas, por considerarmos que sua utilização é muito recente e os participantes ainda não sistematizaram essas formas em seu uso da língua. 6.3.2 Componente de Conhecimento Específico Na parte da prova relativa às questões discursivas do Componente de Conhecimento Específico ( Tabela 6.17 ), observa -se que a Média foi bem mais baixa do que para as questões discursivas do Componente de Formação Geral.
Enquanto no Componente de Formação Geral, a Média para estudantes de Letras-Português e Espanhol ( Licenciatura ) de todo o Brasil foi 56,5, na parte de Conhecimento Específico, a Média foi 31,0.
A maior Média deste componente foi obtida pelos estudantes da região Sudeste ( 37,3 ), e a menor, pelos da região Norte ( 25,0 ).
Quanto à variabilidade das notas, o Desvio padrão de todo o Brasil foi 23,4.
O maior Desvio padrão foi encontrado na região Sudeste ( 25,5 ), e o menor, na região Norte ( 19,7 ).
A maior nota Máxima, 100,0, foi obtida nas regiões Nordeste e Sudeste.
Nas demais regiões a nota Máxima foi : 91,7 na região Norte, 93,3 na região Sul e 91,7 na região Centro-Oeste.
A nota Mínima ( 0,0 ) foi obtida por, pelo menos, um aluno em todas as regiões do Brasil.
A Mediana do Brasil, como um todo, foi 28,3.
A maior Mediana foi obtida nas regiões Sudeste e Sul ( 36,7 ), e a menor, na região Nordeste ( 20,0 ).
No Gráfico 6.10, é representada a distribuição das notas nas questões discursivas no Componente de Conhecimento Específico.
A moda dessa distribuição ocorre no intervalo [ 0 ; 10 ], com quase 25 \% do total de participantes.
Observa -se que nesse intervalo também estão computadas as ocorrências de respostas em branco para as três questões discursivas de Conhecimentos Específicos.
A análise de cada uma destas questões será feita a seguir. 6.3.2.1 Análise da Questão Discursiva 3 do Componente de Conhecimento Específico Na questão 3, cujos resultados aferidos encontram -se descritos na Tabela 6.18, a Média dos estudantes de todo o Brasil foi 30,2, questão de desempenho intermediário dentre as três discursivas de Conhecimento Específico.
A menor Média nessa questão foi obtida pelos alunos da região Nordeste ( 23,6 ), enquanto a maior Média foi obtida na região Sudeste ( 38,7 ).
Quanto à variabilidade das notas, o Desvio padrão de todo o Brasil foi 34,7.
O maior Desvio padrão foi obtido na região Sudeste ( 37,2 ), enquanto o menor foi obtido na região Norte ( 31,8 ).
A nota Máxima, 100,0 pontos, foi alcançada por, pelo menos, um aluno de todas as regiões.
A Mediana do Brasil, como um todo, foi 0,0, a mesma das regiões Norte, Nordeste e Centro-Oeste, indicando que pelo menos 50 \% dos participantes de todo o Brasil e dessas regiões obtiveram nota zero.
Já nas regiões Sudeste e Sul, a Mediana foi 50,0.
A nota Mínima ( 0,0 ) foi a mesma em todas as regiões do Brasil.
No Gráfico 6.11, mostra -se a distribuição das notas na questão discursiva 3, do Componente de Conhecimento Específico da área de Letras-Português e Espanhol ( Licenciatura ).
Sem se levar em consideração que mais de 15 \% dos estudantes deixaram a resposta dessa questão em branco e cerca de 33 \% receberam nota zero, a moda dessa distribuição seria o intervalo ( 40 ; 50 ]. 6.3.2.2 Comentários sobre as respostas à Questão Discursiva 3 Na Questão 3, usou -se como contexto o título de um documentário sobre o cartunista Laerte Coutinho :"Laerte -se".
Além disso, trazia uma breve descrição da temática abordada no documentário, a trajetória de mudança de gênero e os preconceitos contra transgêneros.
Cabe observar que o contexto dessa questão estava fortemente associado à temática explorada por uma das duas questões de Formação Geral, a Questão 2 da mesma prova.
Assim, de certa forma, os enunciados das duas questões influenciaram as discussões apresentadas pelos estudantes em suas respostas.
Dentre os dois itens solicitados, o mais específico, para um formando na área de Letras, era o item ‘ a ’, que solicitava uma explicação do uso da partícula “ se ”, em “ Laerte -se ”, à luz do conceito de inovação linguística, comparando com seu uso como pronome conectado a verbos.
Já a solicitação do item ‘ b ’ não era uma questão específica da área de formação, já que se espera que a capacidade de interpretação de uma expressão escolhida como título de um documentário seja de domínio público.
A maioria dos estudantes começou a resposta pelo item ‘ b ’.
Talvez isso se explique pela sequência das questões discursivas da prova, caso o respondente, motivado pela Questão 2 de Formação Geral, tenha seguido para a Questão 3, de conhecimentos específicos.
Além disso, o item era fácil, envolvendo apenas uma interpretação simples diante da ajuda das informações do enunciado e do fato de se tratar de tema polêmico e bastante atual.
Mesmo assim, ao invés de analisar e interpretar, alguns tenderam a expor suas opiniões pessoais sobre o assunto.
Quanto ao item ‘ a ’, a correta explicação dos verbos reflexivos é um tema que o futuro professor deve conhecer, mas não foi solicitada com base no conceito de inovação linguística e não se exigia uma explicação gramatical, fato este que, além de ser explicitado pelo comando, o padrão de resposta deixou claro, por mostrar que se esperava uma breve explicação de inovação linguística, mais do que uma diferenciação entre os verbos reflexivos e a expressão título do documentário.
A Questão 3 teve desempenho baixo.
Pode -se considerar que teve, para os estudantes, um nível de dificuldade intermediário dentre as três questões discursivas de conhecimentos específicos.
Em geral, os estudantes não utilizaram uma linguagem acadêmica, e muitas respostas tinham um tom opinativo.
Categorias e conceitos gramaticais esperados numa resposta de futuros professores estiveram praticamente ausentes.
A comparação do uso do pronome reflexivo “ se ”, pedida no comando do item ‘ a ’, não foi explorada, as respostas limitavam -se a falar sobre um ou sobre o outro tipo de uso.
O item ‘ a ’ foi sempre respondido de forma muito sintética e com uma linguagem bem coloquial.
Como já mencionado, a maioria deixou o item ‘ a ’ para o final do texto da resposta ou não respondeu a ele, preferindo dissertar sobre o item ‘ b ’.
No item ‘ a ’, pedia -se aos estudantes análise, reflexão e elaboração próprias e não repetição de conteúdo.
A falta de respostas mostra que os estudantes não estavam preparados para situações nas quais eles precisem refletir sobre um conteúdo novo e dar uma resposta original.
As respostas que receberam as notas mais baixas eram curtas, sem que os alunos tivessem feito a análise solicitada no item ‘ a ’ e/ou em que os estudantes se restringiam a expressar opiniões pessoais sobre a questão dos transgêneros, sobre preconceito ou sobre o Laerte.
Nas respostas medianas, geralmente, o item ‘ b ’ foi respondido a contento.
O conteúdo mais específico da formação em Letras, presente na questão e importante na formação para o magistério, os verbos reflexivos, não teve muito peso na questão já que precisava ser abordado apenas com base na inovação linguística.
A questão de interpretação, que deveria ter sido discutida a partir do item ‘ a ’, foi respondida a partir da apresentação da 6.3.2.3 Análise da Questão Discursiva 4 do Componente de Conhecimento Específico Na Tabela 6.19, estão contidas as informações relativas à questão 4 do conjunto de questões do Componente de Conhecimento Específico.
O desempenho dos estudantes nessa questão foi inferior ao das questões 3 e 5.
A Média geral do Brasil foi 25,6, sendo a menor Média registrada na região Centro-Oeste ( 18,9 ), e a maior, na região Sul ( 35,5 ).
A nota Máxima ( 100,0 ) foi atingida por, pelo menos, um concluinte de todas as regiões.
A Mediana em todo o Brasil foi 20,0, a mesma da região Sudeste.
Nas regiões Norte, Nordeste e Centro-Oeste as medianas foram iguais a 10,0, e na região Sul foi igual a 40,0.
A nota Mínima foi zero em todas as regiões.
No Gráfico 6.12, encontra -se a distribuição de notas da questão discursiva 4, do Componente de Conhecimento Específico.
Essa distribuição tem moda na classe de estudantes que não responderam à questão, seguida daqueles que obtiveram notas no intervalo ( 50 ; 60 ], um máximo local, e dos que, tendo respondido à questão, tiraram nota zero. 6.3.2.4 Comentários sobre as respostas à Questão Discursiva 4 Na Questão 4, exigiu -se uma resposta em língua espanhola sobre a análise de um cartaz, que usava linguagens verbal e não verbal, e sobre as vantagens de se utilizarem vários gêneros textuais.
Foi, sem sombra de dúvida, a questão mais difícil das três discursivas de conhecimentos específicos.
O primeiro item era mais prático, e o segundo, mais teórico.
No item ‘ a ’, solicitava -se uma análise do texto ( cartaz apresentado no enunciado ), explicando como as inter-relações das duas linguagens verbal e não-verbal contribuem para a construção de sentidos.
No item ‘ b ’, pediam -se duas vantagens do uso de diferentes gêneros textuais no ensino da língua espanhola.
Cabe destacar que a questão era muito aberta, dando margem a diversas possibilidades de explicação no primeiro item e a uma vasta gama de vantagens que podem ser imaginadas para o uso de diferentes gêneros textuais no ensino de línguas.
O padrão apresentava boa diversidade de respostas possíveis para ambos os itens.
Apesar de no enunciado constar um cartaz sobre o Dia Mundial da AIDS e, abaixo dele, no comando, haver a afirmação “ tendo em conta o cartaz, faça o que pede a continuação", infelizmente, na grande maioria das respostas ao item ‘ a ’, o cartaz foi ignorado.
Um número grande de respondentes tentou explicar como as inter-relações contribuíram na construção do sentido do texto sem que se levasse em conta o cartaz.
Outros simplesmente afirmaram que a inter-relação era muito importante, mas sem qualquer análise.
Talvez o verbo"explicar"do comando do item ‘ a ’ tenha levado o estudante a achar que tinha que defender/justificar uma posição ou dar explicações em abstrato.
A maioria não analisou o cartaz, como se esperava pelo padrão de resposta.
O enunciado do item ‘ b ’ estava claro e não deu lugar a segundas interpretações nem a nenhum tipo de ambiguidade.
No entanto, um número expressivo de estudantes deu exemplos de gêneros textuais, sem que isso lhes tenha sido pedido, e alguns não mencionaram as vantagens.
Na Questão 4, pretendia -se avaliar a capacidade interpretativa do estudante e seu conhecimento teórico de uma metodologia muito utilizada em aulas de língua estrangeira como a utilização e a análise de material variado.
No entanto, como já comentado, poucas respostas conseguiram atender ao que se pedia no enunciado do item ‘ a ’.
O cartaz não foi analisado.
Os estudantes fizeram uma abordagem teórica de um tema que deveria ter sido prática.
No item ‘ b ’, surpreendeu o fato de haver respostas que mencionavam diversos gêneros textuais, sem citar vantagens de seu uso para o ensino-aprendizagem.
Foram muitas as respostas em branco.
Dentre os que responderam, em geral, os estudantes escreveram um espanhol de nível razoável para muito bom.
No entanto, mesmo com um correto uso da língua espanhola, poucos conseguiram chegar a notas mais altas, seja pela desatenção, pelo mal-entendido do enunciado ou, ainda, por desconhecimento.
Poucas respostas cumpriram o esperado pelo padrão.
Considera -se o tema dos gêneros textuais muito pertinente, e a questão buscava avaliar os aspectos prático e teórico do assunto.
A partir de um exemplo concreto, sem teorizar em abstrato, como muitos fizeram, a questão buscava avaliar a capacidade analítica do futuro professor.
Apesar disso, os resultados tornaram difícil precisar o que, de fato, aconteceu para justificar o desempenho tão baixo, muito abaixo do que seria desejável. 6.3.2.5 Análise da Questão Discursiva 5 do Componente de Conhecimento Específico Na Tabela 6.20, \textbf{contam} as informações relativas à questão 5, do conjunto do Componente de Conhecimento Específico.
O desempenho dos estudantes de todo o Brasil nesta questão foi o mais alto dentre as três questões discursivas desse componente.
A nota Média dos estudantes de todo o Brasil foi 37,0.
A maior Média foi registrada na região Sudeste ( 41,9 ), enquanto a menor Média foi registrada na região Norte ( 29,6 ).
Quanto à variabilidade das notas, o Desvio padrão dos alunos do Brasil, como um todo, foi 28,9.
Enquanto o maior desvio foi encontrado na região Sudeste ( 30,7 ), o menor foi encontrado na região Norte ( 22,0 ).
A Mediana para o Brasil e das regiões Norte, Nordeste e Centro-Oeste, foi 25,0.
Nas regiões Sudeste e Sul, a Mediana foi maior, igual a 50,0, em ambas.
Para o conjunto de alunos de Letras-Português e Espanhol ( Licenciatura ) do Brasil, a nota Máxima foi 100,0, a mesma de quatro das regiões.
A exceção foi a região Norte, onde a nota Máxima foi 75,0.
A nota Mínima ( 0,0 ) foi a mesma em todas as regiões.
Sudeste ( 41,9 ), enquanto a menor Média foi registrada na região Norte ( 29,6 ).
Quanto à variabilidade das notas, o Desvio padrão dos alunos do Brasil, como um todo, foi 28,9.
Enquanto o maior desvio foi encontrado na região Sudeste ( 30,7 ), o menor foi encontrado na região Norte ( 22,0 ).
A Mediana para o Brasil e das regiões Norte, Nordeste e Centro-Oeste, foi 25,0.
Nas regiões Sudeste e Sul, a Mediana foi maior, igual a 50,0, em ambas.
Para o conjunto de alunos de Letras-Português e Espanhol ( Licenciatura ) do Brasil, a nota Máxima foi 100,0, a mesma de quatro das regiões.
A exceção foi a região Norte, onde a nota Máxima foi 75,0.
A nota Mínima ( 0,0 ) foi a mesma em todas as regiões.
No Gráfico 6.13, representa -se a distribuição das notas da questão discursiva 5 do Componente de Conhecimento Específico.
Nessa distribuição, diferentemente das outras duas questões discursivas, apesar do grande número de respostas em branco, mais do que 20 \% dos participantes, a moda principal ocorreu no intervalo ( 20 ; 30 ].
Tal resultado está associado à grande frequência de notas 25, devido à anulação de parte do item ‘ b ’ da questão. 6.3.2.6 Comentários sobre as respostas à Questão Discursiva 5 Na Questão 5, exigia -se do aluno conhecimento em literatura em geral e em literatura brasileira em particular.
Do enunciado constava um trecho do Auto da compadecida de Ariano Suasuna e, considerando o texto para responder aos itens da questão, o estudante deveria conhecer as convenções de um “ auto ”, saber identificar uma figura da linguagem e explicar o efeito que essa expressão tem no texto.
O enunciado era claro, menos a segunda parte do item ‘ b ’ ( o efeito no texto ) que precisou ser anulada, já que o fragmento apresentado finalizava com a expressão em questão e não era possível, de forma alguma, que o aluno soubesse o efeito que ela teria no texto, já que não havia mais texto.
Caso a questão tivesse solicitado o efeito da figura de linguagem no leitor, teria sido possível responder ao item.
Se fosse esse o caso, faria sentido que o texto acabasse com a expressão.
E, se fosse esse o caso, a Questão 5 avaliaria a mesma competência solicitada no item ‘ b ’ da Questão 3, o efeito do texto no interlocutor.
O resultado geral foi um pouco melhor do que o das duas outras questões discursivas.
O fragmento escolhido pertence a um texto importante na história da literatura brasileira, portanto, provavelmente era conhecido pela maioria dos estudantes.
No item ‘ a ’, apenas se pedia que o estudante indicasse duas convenções, e o padrão de resposta ofereceu cinco possibilidades de resposta.
No item ‘ b ’, a exigência foi um pouco maior.
Nele, o estudante deveria acertar o nome da figura de linguagem.
Houve bastante diversidade nas respostas dadas ao item ‘ a ’ e, apesar de não ser solicitado, muitas vezes, o estudante exemplificava com o próprio texto.
No item ‘ b ’, os estudantes erraram muito o nome da figura de linguagem e também na interpretação da expressão, apesar de terem respondido ao que parecia mais natural, o efeito da figura de linguagem no leitor.
A linguagem utilizada nas respostas foi adequada, mas, frequentemente, as respostas eram opinativas.
O tema do racismo deu margem a que os estudantes fugissem do texto literário e defendessem ideias, especialmente em relação à parte que foi anulada, na interpretação do eufemismo.
No entanto, também foi bastante comum o estudante tentar inferir as convenções do subgênero “ auto ” a partir do próprio texto.
Como a segunda parte do item ‘ b ’ foi anulada e houve um predomínio de respostas com notas baixas e medianas, depreende -se que o desempenho real dos estudantes foi insatisfatório.
Esperava -se que o nível das respostas fosse melhor em se tratando de temática literária.
Os estudantes tinham de identificar apenas duas convenções do texto e, mesmo assim, tiveram dificuldades em expressá -las.
Poucos conheciam o nome da figura de linguagem em questão e, por isso, descreviam a figura sem nomeá -la.
Nas respostas dadas pelos estudantes à Questão 5 observou -se que o tema do racismo foi o principal foco, ainda que não fizesse parte diretamente do enunciado da questão.
Mais uma vez, os estudantes fugiram do que foi solicitado na questão, redigindo textos com opiniões pessoais, que não se configuravam como resposta. 6.3.3 Considerações Finais As três questões pautaram -se em contextos atuais, importantes, porém polêmicos e um tanto distantes de aspectos necessariamente presentes nos currículos dos cursos e nas Diretrizes Curriculares.
Obviamente, a formação superior precisa ser ampla e comprometida com as questões sociais, culturais e ambientais, bem como deve desenvolver competências para ler, interpretar e produzir textos com clareza e coerência.
No entanto, essas competências gerais, dentre várias outras, necessárias à formação dos cursos superiores, já são avaliadas pelas questões de Formação Geral.
Assim, as questões do componente específico poderiam envolver situações-problema e estudos de caso voltados para Nessa edição do Enade, as questões discursivas de conhecimentos específicos acabaram sendo “ contaminadas ” pela forma com que os estudantes costumam responder às questões de Formação Geral.
De fato, especialmente as questões 3 e 4 foram construídas a partir de temáticas que também foram foco das questões 1 e 2, respectivamente doenças sexualmente transmissíveis e identidade de pessoas trangêneras.
O efeito dessa similaridade temática acabou se observando também na Questão 5, na qual diversos estudantes focaram a discussão do racismo.
Os itens discursivos abrangeram diferentes áreas do curso de Letras Português-Espanhol : Linguística, Língua Espanhola, Literatura Brasileira, Metodologias de Ensino de Língua Espanhola.
Quanto aos conteúdos específicos da área de Letras Português e Espanhol avaliados, destacaram -se a capacidade de interpretação de textos, presente nas três questões, o conceito de inovação linguística ( Questão 3 ), os conceitos de linguagens verbal e não verbal ( Questão 4 ), características do gênero “ auto ” e figuras de linguagem ( Questão 5 ).
Do ponto de vista dos conhecimentos pedagógicos, necessários ao exercício da docência, avaliou -se o reconhecimento da importância do uso de gêneros textuais variados ( Questão 4 ), e a formação em Espanhol pôde ser observada na redação da resposta à Questão 4.

% matched lemmas: contar, panorama
\end{textsample}
